% !TeX program = xelatex
% 编译：在 docs 目录执行两次 xelatex recast-project-description.zh-CN.tex。
% 依据项目技术方案及 2026-09-08 实现说明编写；第 3、8、9、10 章正文留空。
\documentclass[UTF8,12pt,openany,fontset=fandol]{ctexrep}

\usepackage[a4paper,top=2.5cm,bottom=2.5cm,left=2.7cm,right=2.7cm]{geometry}
\usepackage{amsmath,amssymb}
\usepackage{booktabs,longtable,array,tabularx}
\usepackage{enumitem}
\usepackage{xcolor}
\usepackage{hyperref}

\definecolor{ProjectBlue}{HTML}{2255A4}
\hypersetup{
  colorlinks=true,
  linkcolor=ProjectBlue,
  urlcolor=ProjectBlue,
  citecolor=ProjectBlue,
  pdfauthor={RECAST 项目团队},
  pdftitle={RECAST：面向大型工业软件的智能检索与自适应开发平台},
  bookmarksnumbered=true
}
\setlength{\parindent}{2em}
\setlength{\parskip}{0.2em}
\setlength{\emergencystretch}{2em}
\renewcommand{\arraystretch}{1.3}
\setcounter{tocdepth}{2}
\setcounter{secnumdepth}{2}
\setlist[itemize]{leftmargin=2.2em,itemsep=0.3em,topsep=0.3em}
\setlist[enumerate]{leftmargin=2.4em,itemsep=0.3em,topsep=0.3em}
\newcolumntype{P}[1]{>{\raggedright\arraybackslash}p{#1}}
\newcolumntype{Y}{>{\raggedright\arraybackslash}X}

\begin{document}

\begin{titlepage}
  \centering
  \vspace*{2.4cm}
  {\zihao{1}\heiti RECAST\par}
  \vspace{0.6cm}
  {\zihao{2}\heiti 面向大型工业软件的\\[0.35em]
  智能检索与自适应开发平台\par}
  \vspace{1.6cm}
  {\zihao{2}\heiti 项目说明书\par}
  \vfill
  {\zihao{4}RECAST 项目团队\par}
  \vspace{0.6cm}
  {\zihao{4}2026 年 9 月\par}
  \vspace*{1.5cm}
\end{titlepage}

\pagenumbering{Roman}
\tableofcontents
\clearpage
\pagenumbering{arabic}

% 第一章：产品目标与定位。
\chapter{项目概述}
\label{chap:overview}

\section{项目简介}

RECAST 是一套面向大型工业软件与企业代码资产的智能检索与自适应开发平台。用户通过自然语言描述需要开发、修改或迁移的功能，系统结合源码解析、模块建模和混合检索，定位相关实现及依赖，并组织成包含源码位置、版本和证据的任务上下文，供开发人员与 Coding Agent 使用。

平台围绕两条工作流建设。任务检索服务于大型工程中的功能定位和代码理解，向 Codex、Claude Code 等编程智能体提供开发依据；复用迁移服务于企业已有代码的适配与再利用，在参考实现检索之后继续开展分析规划、多文件翻译和编译反馈修复。两条工作流共享代码索引、模块表示和检索能力。

项目面向百万至千万行级、多语言工程的建设需求，重点研究依赖约束的自适应分层并行建模、跨语言依赖恢复、多粒度级联检索，以及任务覆盖与预算约束的上下文选择。已有原型完成了百万行级真实源码的索引和任务上下文返回验证，千万行规模与企业任务效果按独立指标继续验证。

\section{项目定位}

RECAST 定位为企业代码资产与 AI Coding 工具之间的工程理解支撑平台。平台承担可重复使用的工程扫描、结构分析、索引维护和基础检索，将模型分析集中到相关模块与有限候选范围。Coding Agent 在取得上下文后负责制定修改方案、生成代码并执行验证。

平台的交付单位是围绕当前任务组织的代码上下文，包含主要实现、必要依赖、来源版本与尚未确认的信息，支持开发者核查和智能体继续查询。模块树表达功能职责与包含关系，依赖图表达调用和工程关联，两者共同支撑实现定位与范围分析。

\section{核心应用场景}

\begin{enumerate}
  \item \textbf{功能定位与需求迭代。} 根据自然语言需求，在当前工程中定位实现入口、关键分支和相关配置，支持已有功能修改与增量开发。
  \item \textbf{企业代码复用与迁移。} 从参考工程中发现可借鉴实现，识别接口、类型及外部依赖，为跨语言、跨框架或跨平台适配提供依据。
  \item \textbf{跨模块影响分析。} 围绕公共接口、基础类型或配置变更，查找相关调用方、实现模块及验证依据，辅助确定修改范围。
  \item \textbf{多语言工程理解。} 结合语言内分析与桥接关系恢复，追踪上层接口、注册入口和底层实现之间的关联。
  \item \textbf{Coding Agent 上下文供给。} 将检索结果整理为预算内的源码与依赖证据，通过 HTTP、MCP 和 IDE 接入下游开发流程。
\end{enumerate}

上述场景主要面向维护大型、多模块、多语言软件的企业研发团队。功能定位是统一入口，复用迁移与影响分析在相同代码模型上选择不同范围和关系，减少重复建库。

\section{建设目标与核心能力}

项目以大型工程的可处理性、任务定位质量和下游开发效果为共同目标，形成表~\ref{tab:capability-goals} 所示的能力体系。

\begin{table}[htbp]
  \centering
  \caption{项目建设目标与核心能力}
  \label{tab:capability-goals}
  \small
  \begin{tabularx}{\textwidth}{P{0.18\textwidth}YY}
    \toprule
    建设目标 & 核心能力 & 面向用户的结果 \\
    \midrule
    大仓建模 & 分批解析、分层模块理解与增量索引 & 可浏览、可查询的工程结构 \\
    依赖恢复 & 编译语义与跨语言注册绑定分析 & 可追溯的调用与工程关系 \\
    任务定位 & 多粒度混合召回与候选排序 & 与需求相关的实现位置 \\
    上下文构建 & 必要证据选择、去重与预算控制 & 可供 Agent 使用的源码上下文 \\
    复用迁移 & 参考检索、分析规划与生成验证 & 可审阅的代码适配结果 \\
    \bottomrule
  \end{tabularx}
\end{table}

检索准确率、响应时间和代码功能完备度采用任务集与统一环境评价。模块数量、文件覆盖率及索引条数说明工程处理情况，任务定位和功能验证结果则反映业务效果。

\section{术语与业务角色}

平台统一使用“目标工程”表示当前需要开发、修改或适配的工程，“参考工程”表示提供可借鉴实现的工程。两者是任务中的角色，同一个工程可在不同任务中承担不同角色。检索返回的模块称为“候选模块”，用户选定后称为“参考模块”。

“代码仓库”用于描述源码来源和版本管理，“项目”用于描述构建单元，“模块”用于描述功能职责。一个项目可包含多个模块，模块之间通过依赖关系协作。文件、类、函数与代码块构成更细粒度的检索和源码读取单位。

本文将交付给 Agent 的源码、关系及来源信息统称为“任务上下文”，英文简称 Context；其结构化输出称为上下文包。上下文预算指允许交付的 token、文件数和源码行数等上限。

\section{当前实现范围}

当前原型采用 Node.js 与 TypeScript 实现核心服务，使用 SeekDB 保存代码结构和检索投影，通过浏览器工作台、VS Code 扩展及 HTTP/MCP 提供访问。已接通源码快照、分批解析、模块建模、混合检索、有限依赖补充和真实源码上下文返回。

模块建模已加入父子层级、自适应停止、聚合统计与层次展示，并保留结构判断和可选模型判断入口。当前真实层次验证采用本机结构基线，模型职责判断的效果另行评价。完整跨语言图谱、持久分支并行、任务感知精排和覆盖驱动的上下文选择属于后续增强范围。第~\ref{chap:implementation} 章说明实际功能及其接口边界。

% 第二章：需求依据与业务问题。
\chapter{背景与需求分析}
\label{chap:requirements}

\section{行业背景与工程需求}

\subsection{存量代码资产与软件持续演进}

大型工业软件在长期研发过程中积累了大量业务实现、平台适配代码和公共基础组件。这些内容通常分布于不同仓库、技术栈和历史版本中，研发人员需要投入较多时间理解模块职责、确认已有能力并追踪修改影响。平台升级和国产软件生态适配进一步增加了对存量代码理解与复用的需求。

企业复用代码时既要定位相近功能，也要确认接口、运行环境和依赖是否适合目标工程。仅凭模块名称或局部代码相似性，难以判断完整适配范围。由代码结构、依赖与任务信息共同支撑的检索能力，是复用迁移的重要基础。

\subsection{AI Coding 的工程上下文需求}

Coding Agent 可以调用工具搜索、读取和修改代码。面对陌生工程，Agent 需要逐步确定实现入口、关联类型、配置与调用链，然后形成修改方案。输入上下文缺少关键依赖时，生成结果可能遗漏行为；输入包含大量无关内容时，又会增加阅读成本并干扰判断。

大型工程难以作为一次模型输入整体提供。系统需要将完整工程转化为可持续查询的索引，再根据任务选择源码与依赖，使模型的输入范围与实际修改范围相匹配。

\subsection{多语言与隐式依赖的复杂性}

工业软件中的功能往往跨越业务模块、基础库和平台接口。语言内关系可通过导入、调用和继承表达，跨语言关系还可能经过 Native 注册、接口描述或构建配置连接。相同名称在不同构建单元和依赖版本中也可能指向不同实现。

因此，源码文本、语法结构和构建信息需要协同分析。关系解析既要记录成功连接的端点，也要保留未解析和存在歧义的位置，以便任务检索与人工核查识别信息缺口。

\section{企业需求与项目响应}

项目同时承接两类相关需求：华为场景关注 AI Coding 超大工程代码理解与快速索引，潍柴场景关注企业代码复用迁移。前者要求根据开发任务定位相关代码并交付上下文，后者在此基础上继续分析与翻译参考实现。

两类需求共享工程接入、索引和检索底座，区别主要在输出之后的工作流。任务检索可独立提供给外部 Agent；复用迁移进一步引入目标接口约束、名称映射、代码生成和编译反馈。通过共享索引与证据，减少两条业务链路中重复的代码读取与分析。

\section{命题目标与验收口径}

\subsection{目标语言与性能要求}

依据项目留存的企业需求，目标语言包括 Java、ArkTS、Kotlin Multiplatform 和 C/C++。检索准确率目标至少为 90\%，平均检索耗时在 10 秒以内，生成或转换代码的功能完备度目标超过 90\%。百万至千万行规模是建库与查询能力的验证范围。

这些数值是建设与验收目标。不同语言的语法解析、编译语义和跨语言绑定能力分别验证，语言支持范围以具体工程和构建场景为依据。

\subsection{从业务指标到可测量任务}

检索评价需要固定任务、有效实现位置和必要证据。定位质量可通过前若干候选是否包含有效入口、相关实现的召回比例和排序位置衡量。对于不存在匹配实现的任务，还应评价正确返回无匹配的能力。

查询时延从用户或 Agent 发起请求开始，计入排队、编码、召回、重排、依赖扩展、源码读取和上下文组装。功能完备度以预先确定的行为要求和测试断言衡量，编译通过单列为工程检查结果。

\section{典型业务任务}

\subsection{自然语言功能定位}

用户提出“修改文件上传总大小和单文件大小限制”时，系统需要找到请求解析流程、大小阈值及对应检查分支。代码中的标识可能是 \texttt{sizeMax}、\texttt{fileSizeMax} 或具体异常类型，与用户表述存在差异。检索既要利用自然语言语义，也要保留接口名和代码标识的精确匹配。

Commons FileUpload 是当前可用于验证这一任务的真实语料。任务结果需要给出对应实现和源码范围，支持用户区分实际限制检查与仅在注释、日志或说明中提及上传的内容。

\subsection{跨平台功能迁移}

迁移上传功能时，需求会从“找到限制检查”扩展到“保留限制、异常处理和资源释放行为”。相关实现可能依赖流接口、文件项工厂、平台配置及测试。系统需要明确源实现、目标工程已有接口和需要适配的外部依赖，再将这些约束交给生成阶段。

涉及 ArkTS 与 Native 层的迁移还需识别注册端点和底层实现。该场景用于指导跨语言关系设计，具体效果需在对应可构建工程上验证。

\subsection{公共接口修改与影响分析}

修改公共接口时，除接口定义外，还需定位调用方、实现类及受影响验证范围。此类任务强调从被修改对象向使用者扩展；实现迁移则更关注被调用对象和运行依赖。不同任务对同一代码入口需要不同方向的关系查询。

\section{现有检索方式的适用范围}

关键词与全文检索适合已知名称、错误码和固定字符串，具有直接的位置依据。符号导航在具备语义索引时支持定义和引用追踪。向量检索可以连接自然语言描述与功能相近的代码表示。Agent 多轮探索则能够结合新发现调整搜索方向。

这些方式在系统中承担互补职责。全文与符号检索提供精确候选，语义检索扩大描述不一致时的召回范围，工程关系补充实现之间的关联，Agent 在候选范围内完成进一步判断。设计重点是确定各方法的输入、组合顺序和计算预算。

\section{需求归纳}

项目需求可归纳为四项：在受控资源内组织并维护大型代码仓；恢复有源码和构建依据的跨语言关系；在有限时延内准确定位任务实现；将必要证据组织为低冗余、可核查的上下文。四项需求依次对应建模、关系、检索和上下文技术，同时通过统一的版本与范围贯穿完整流程。

% 第三章按要求保留结构，正文留空。
\chapter{市场现状与竞品分析}
\section{市场供给现状}
\section{主要产品对比}
\section{竞争判断与项目差异化}
\section{竞品对标方法}

% 第四章仅定义挑战和设计目标，方法集中在第六章。
\chapter{技术挑战与设计目标}
\label{chap:challenges}

\section{超大规模代码的组织与持续维护}

\subsection{全量分析的资源与更新开销}

源码扫描、语法解析、符号绑定、模型理解和向量生成具有不同的计算与存储需求。将各阶段组织为一次全量作业，会使内存驻留、数据库写入和模型等待相互叠加。局部修改之后重复执行整仓分析，也会增加不必要的重算。

持续维护的难点在于确定变化的影响范围。文件正文、公开接口和构建配置的变化具有不同传播范围；仅重建改动文件可能遗留旧关系，直接重建全仓又会损失增量收益。系统需要同时考虑输入变化、结果依赖和查询可见性。

\subsection{结构粒度与功能职责的不一致}

目录与包提供了自然的工程组织线索，但不能直接等同于功能模块。同一目录可能包含多种职责，同一功能也可能跨越多个目录。模块数量多或文件归属完整，并不意味着边界符合实际职责。

不同规模和复杂度的工程也不适合固定层数。小工程可能只需要一个或几个叶模块，大型业务分支则需要继续细分。因此，模块树应表达职责层次，并允许不同分支在不同深度终止。

\subsection{设计目标}

建模过程需要控制解析批次与并行任务的资源开销，允许基础索引与增强理解分阶段交付。模块结果应满足文件归属明确、父子范围一致和职责证据可追溯，变更后优先更新受影响结果。评价同时关注资源、构建时间、职责边界与增量重算范围。

\section{跨语言关系的恢复与可信使用}

\subsection{语言内符号歧义}

同名声明、继承、泛型和条件编译可能使语法引用对应多个目标。调用是否可用还取决于依赖版本、源集和编译选项。系统需要区分语法层观察到的引用与在特定构建条件下确认的绑定，避免将所有候选都当作实际调用。

\subsection{语言间桥接与隐式关联}

Native 注册、接口描述和框架配置将不同语言实体连接起来，关系证据分散在多个文件中。恢复链路需要识别注册键、签名和作用域，单纯匹配名称或向量相似度无法确定真实绑定。注册信息变化还可能使先前推导的关系失效。

\subsection{设计目标}

关系模型需要统一表达不同语言的声明、调用和绑定，保存推导依据及适用的构建范围。查询时区分已解析、歧义和未解析关系，更新时撤回依赖旧证据的结果。评价按语言、关系类型与构建场景分别统计，并检查关键路径是否能够恢复。

\section{检索质量与在线时延的平衡}

\subsection{多粒度候选与信号差异}

函数修改、模块复用和子系统理解需要不同粒度的检索结果。若只在全仓细粒度索引中查找，候选噪声可能较多；若只沿模块树下钻，则可能因模块边界或路由错误遗漏相关实现。全文、向量和模块视图的排名又具有不同含义，需要合理融合。

\subsection{精排与图扩展的计算成本}

模型重排能够比较复杂任务条件，但候选过多会提高输入量与响应时间。依赖扩展有助于补充上下文，高扇出节点却可能迅速扩大查询规模。系统需要把较昂贵的处理限制在有价值的候选范围内，并将各阶段计入统一时延预算。

\subsection{设计目标}

检索应同时利用层次结构与全局召回，保留精确标识和语义描述的互补优势。用户指定的粒度和工程范围必须得到执行。评价结合有效入口命中、实现召回、排序质量、无匹配任务表现和完整上下文交付时延。

\section{任务证据覆盖与上下文预算}

\subsection{相关结果与必要信息的差异}

相似代码不一定包含完成任务所需的全部信息。修改接口可能需要调用方约束，迁移实现可能需要构建配置和异常行为。另一方面，同一函数可能通过声明、源码块和模块摘要多次命中，重复内容会消耗上下文预算。

\subsection{任务覆盖的判定难点}

必要证据并不是全图中所有可达节点。系统需要根据任务区分核心实现、接口约束与支持性内容，并判断缺少哪些信息会影响开发。在线判断依据已解析源码和任务条件，动态行为或缺失构建信息可能使部分依赖仍不确定。

\subsection{设计目标}

上下文选择以必要证据覆盖为目标，以 token、文件数和源码行数为约束。输出保留核心源码及来源，对未覆盖内容提供明确线索。方法效果通过必要证据覆盖和 Agent 任务完成情况共同评价，在此基础上衡量总 token 与额外查询成本的变化。

% 第五章：系统组成、数据流及部署分工。
\chapter{产品总体方案}
\label{chap:architecture}

\section{总体设计思路}

RECAST 采用“离线建模、在线检索、按需交付”的总体设计。离线侧将工程源码转换为结构事实、模块表示和依赖关系，形成可重复使用的索引；在线侧根据当前任务定位实现，补充必要关联，并组织预算内上下文。代码生成和编译验证通过下游 Agent 或复用迁移服务完成。

系统围绕三个共同约束组织各模块。首先，源码与分析结果绑定明确版本，查询使用一致的读取范围。其次，构建和检索分别控制输入量、并行度及执行预算，避免局部查询触发整仓分析。最后，模块说明与依赖关系保留源码依据，最终上下文能够定位到具体文件和实现范围。

任务检索与复用迁移共享底层工程身份和索引。检索范围由本次任务选择，目标写入范围由开发宿主管理，因此参考工程既可参与搜索，也可作为翻译来源，而不会因参与检索自动成为修改目标。

\section{系统总体架构}

系统分为工程接入、代码分析、索引存储、任务检索和应用交付五层。各层通过统一的数据契约交互，核心职责如表~\ref{tab:architecture} 所示。

\begin{table}[htbp]
  \centering
  \caption{系统分层与模块职责}
  \label{tab:architecture}
  \small
  \begin{tabularx}{\textwidth}{P{0.18\textwidth}YY}
    \toprule
    层次 & 主要组成 & 核心职责 \\
    \midrule
    工程接入层 & 仓库登记、项目发现、源码快照 & 确定工程身份、构建范围与输入版本 \\
    代码分析层 & 结构解析、语义分析、模块建模 & 提取符号与关系，组织功能职责 \\
    索引存储层 & 结构事实、模块产物、全文与向量索引 & 持久化分析结果，提供局部查询 \\
    任务检索层 & 查询组织、混合召回、重排、上下文选择 & 定位实现并交付必要证据 \\
    应用交付层 & 工作台、IDE、HTTP/MCP、迁移服务 & 展示结果、接入 Agent 与后续开发流程 \\
    \bottomrule
  \end{tabularx}
\end{table}

模块树与依赖图位于分析和检索之间，承担不同的组织作用。模块树将工程表示为职责层次，为范围收缩、模块选择与摘要检索提供依据。依赖图保留跨文件、跨模块及跨语言关联，为实现之间的连接、影响分析和上下文补充提供依据。向量索引中的近邻图仅用于检索加速，不承担工程关系的含义。

\section{统一代码模型与数据组织}

\subsection{源码事实与功能表示}

基础代码模型包含仓库、版本、项目、文件、符号和源码块。文件记录语言、角色、内容哈希和解析情况；符号记录声明、签名、范围和所属文件；源码块关联具体实现及其位置。构建配置、测试和接口描述作为工程信息纳入相应范围。

模块表示在基础事实之上组织职责、公开接口和边界依赖。叶模块直接拥有文件，父模块汇总子模块的范围与说明。源码仍可按类、函数或代码块独立检索，模块树较浅不影响基础符号的完整索引。

依赖关系记录来源端、目标端、关系类型和解析依据。跨语言增强进一步纳入调用点、注册入口及绑定实体，使同一功能的实现路径可以跨越语言边界。模块摘要用于理解职责，关系事实用于确认代码连接，二者在检索结果中分别呈现。

\subsection{版本与结果发布}

系统使用源码快照固定一次分析输入，并将结构、模块和检索投影关联到相应版本。增强架构允许不同层具有独立构建进度，但对外发布时需要检查它们是否使用兼容输入。一次查询固定其读取视图，后续全局活动版本变化不改变该请求已取得的证据来源。

构建中结果与已发布结果分开管理。新版本完成规定的完整性与可查询性检查后，才切换活动指针；失败作业保留状态和原因，上一有效版本继续提供服务。局部源码读取、模块范围与向量命中必须在同一版本约束下解释。

\section{离线建库与在线检索协同}

\subsection{离线与后台建库}

离线流程首先登记工程、识别项目并捕获源码，然后按受限批次执行结构解析和依赖绑定，将文件、符号及源码块写入索引。基础索引就绪后，系统可提供源码定位；模块理解与增强关系分析继续产生更高层的职责与关联表示。

分层理解以有限范围为分析单元。模型取得候选组、接口和局部证据后决定拆分或停止，宿主检查划分是否合法，再保存节点和子任务。全文与向量表示根据已发布的代码块、叶模块和上层模块分别构建。

\subsection{在线任务处理}

在线请求包含任务描述、工程范围、检索粒度和上下文预算。服务检查所选范围与索引能力，组织声明、实现和模块候选，执行融合排序，读取有限源码并补充相关依赖。最终由上下文构建器输出结构化包和面向 Agent 的文本。

在线路径优先使用按身份、范围和邻接关系的局部查询。扩大任务范围或补充源码通过后续请求完成，不要求为一次查询重新加载完整工程。取消、超时和预算不足在结果中保留相应状态，界面避免使用迟到响应覆盖新的任务。

\subsection{增量维护与反馈}

源码变化后，系统通过内容比较识别可复用和需要更新的结果。增强设计继续沿符号、关系证据和模块归属传播影响，重建相关索引与摘要。涉及接口或配置的变化可能扩展到相邻模块及共同祖先，而正文内部变化通常具有较小影响范围。

编译或行为验证反馈也可成为新的查询输入。Agent 提供错误位置、已修改内容和任务约束，检索侧据此补充相关实现和依赖。连续开发中的未提交修改需要单独的工作树覆盖机制：变化文件替代旧内容，失效关系重新计算，结果明确对应当前读取版本。

\section{两条业务工作流}

\subsection{任务检索工作流}

\begin{enumerate}
  \item 选择目标工程或参考范围，输入开发任务，可附已知模块或代码位置。
  \item 指定主要结果粒度与内容预算，或使用自动模式。
  \item 检索服务召回实现候选，组织关联源码和依赖证据。
  \item 用户查看位置、来源和缺口，预览、选择或导出上下文。
  \item 外部 Agent 通过同一服务取得上下文，并按需继续读取符号和源码。
\end{enumerate}

该流程允许用户在不知道模块位置时直接提出需求。模块导航既可作为浏览入口，也可成为检索范围，不是发起所有任务的前置条件。

\subsection{复用迁移工作流}

复用迁移先明确目标工程的任务与接口约束，再从参考工程中检索候选实现。选定参考方案后，分析阶段整理源码范围、类型与名称映射以及需要处理的外部依赖；生成阶段据此执行多文件适配，并根据编译反馈修复问题。

完整产品流程还需要将生成结果与目标基线、验证记录及差异审阅关联。编译验证用于确认语言和工程集成条件，行为验证用于确认必须保留的功能要求。审阅后的回填属于开发宿主的写入流程，与只读检索接口分别管理。

\section{部署与集成方案}

\subsection{本地化部署}

原型采用本机服务部署，由索引与查询服务访问工程文件和 SeekDB，前端负责交互与展示。向量模型提供查询和文档编码，生成模型用于模块理解或翻译。两类模型承担不同任务，可分别配置；基础结构查询不依赖生成模型。

企业部署可围绕相同服务边界安排索引工作进程、模型服务与数据库资源。资源配置需要分别考虑解析吞吐、模型调用和数据库容量，并保持工程访问范围与模型数据使用方式一致。跨主机部署需要补充相应身份管理和服务配置，当前本机接口不直接等同于分布式部署能力。

\subsection{HTTP、MCP 与 IDE 接入}

HTTP 面向工作台和系统集成，返回结构化上下文包；MCP 面向编程 Agent，提供任务检索与基础代码查询工具；IDE 宿主负责当前工程身份、用户选择、结果展示和开发操作。三种入口复用任务检索服务，减少不同客户端之间的结果差异。

向量模型密钥、数据库连接和工程文件访问由服务或宿主管理。前端传递任务意图和已登记对象身份，检索结果带回版本、源码范围和证据。模型接入与代码写入按照各自工作流配置，避免让检索请求同时承担执行修改的职责。

% 第六章：方法定义、机制及协同。
\chapter{核心技术与项目创新}
\label{chap:methods}

\section{总体技术路线}

项目以“结构与语义协同建模、多粒度检索、任务证据组织”为技术主线。结构分析提供源码和依赖事实，模型在有限证据上识别模块职责；层次结构引导检索范围，全文与向量检索联合召回实现；依赖扩展和上下文选择将候选转化为可供开发使用的证据。

本章阐述完整技术方案及其设计目标。已有结构、索引和查询能力作为基础，分支并行、跨语言绑定和覆盖驱动选择按各自机制逐步增强。实际产品接通情况在第~\ref{chap:implementation} 章集中说明。

\section{依赖约束的自适应分层并行建模}
\label{sec:hierarchical-modeling}

本项技术将大型代码仓的分析分为基础事实构建和分层职责理解两个阶段。基础阶段控制扫描、解析与写入的批次资源；理解阶段利用代码证据决定模块边界和细化深度，通过子任务调度与结果复用降低重复分析。

\subsection{有界解析与分片索引构建}

源码扫描、文件读取、语法解析和索引写入被组织为连续流水线。各阶段按文件数与字节量限制待处理数据，下游达到容量上限时暂停上游生产，避免写入或模型速度不足造成队列持续积累。解析进程按批次回收，使原生解析资源在阶段结束后释放。

基础解析提取文件角色、声明、符号范围和依赖线索。跨文件绑定在必要结构事实可用后完成，已解析关系与未解决引用分别持久化。源码内容采用快照与范围关联，后续检索读取匹配版本的实际文本。

分片构建将单次写入和失败影响限制在局部批次，发布检查仍针对所需范围的完整性。基础索引完成后先提供查询，模块树和增强关系按兼容版本继续构建。当前实现已控制解析批次并采用分批写入，结构元数据仍有全局驻留需求，进一步的分片查询与持久恢复用于继续降低大仓资源开销。

\subsection{结构与语义协同的模块划分}

模块划分以项目、目录和包作为初始候选线索，同时引入公开接口、引用方向和组间依赖统计。候选组将大仓转换为较少的局部分析对象，模型据此识别职责、解释边界，并决定哪些组应合并、哪些范围需要细分。

模型读取的主要内容是候选规模、代表文件、接口、局部源码和依赖证据，而不是每层重复输入完整工程。模块描述需要关联输入证据，避免职责判断脱离实际实现。结构关系为划分提供约束，但高跨组依赖不一定意味着应合并：公共基础库可能具有大量外部消费者，仍可保持独立职责。

文件归属采用叶节点唯一拥有原则。设父模块 $m$ 的文件集合为 $\mathcal{F}(m)$，其有效子模块为 $m_1,\ldots,m_k$，一次完整拆分需要满足：
\begin{equation}
  \bigcup_{i=1}^{k}\mathcal{F}(m_i)=\mathcal{F}(m),
  \qquad
  \mathcal{F}(m_i)\cap\mathcal{F}(m_j)=\varnothing\quad(i\ne j).
  \label{eq:module-partition}
\end{equation}

其中 $k\geq 2$，父节点通过子节点聚合范围，不重复拥有其后代文件。宿主同时检查父节点有效性、树中无环和证据归属。未能有效分析的范围保留明确状态，不能通过丢弃文件满足表面覆盖。

\subsection{自适应细化与分支并行}

分层深度由职责和规模共同决定。一个范围职责单一且已有证据足以解释时，可直接成为叶模块；存在多个可区分职责时，再生成有效子组。代码规模用于提示分析压力，不能单独决定必须拆分。证据不足、无法形成有效边界或达到预算上限时，系统停止当前递归，并记录具体原因。

该机制使不同分支具有不同深度。小工程可以直接形成一个叶模块，大型分支可继续拆解为多个职责单元。模块与子系统是语义类型，层数是结构属性，前端展开深度是展示策略，三者分别管理。

并行调度遵循父子依赖：父节点完成有效划分后，独立子任务可以并行运行；父摘要汇总等待直接子结果，而其他分支可以继续向下细化。解析、模型理解与向量生成采用独立资源池，分别约束并行度和输入预算。

持久分支调度方案使用任务队列传递执行工作，由索引存储保留权威任务状态与产物。任务重试需要校验输入版本和当前执行身份，防止超时结果覆盖新的分析结果。队列任务结束与索引正式发布分别判断，保证部分失败时仍能保留上一有效范围。

\subsection{变更影响传播与局部增量更新}

增量维护首先比较源码内容、分析工具和表示配置，识别可复用输入。正文变化触发相关符号、代码块与模块说明更新；公开接口、构建配置或注册信息变化还需扩展到依赖其内容的关系及模块。摘要向必要祖先传播，直到上层职责和边界影响得到处理。

向量复用依据实际编码内容及模型身份，而非仅依据文件路径。模型、提示策略或表示方式变化时，对应缓存需要失效。未变化的源码和表示可以复用，新结果写入新的可发布版本。

增量正确性通过与相同输入下的全量构建比较验证。结构事实和关系比较规范化后的结果；模型生成的模块描述则比较覆盖、边界、证据及检索效果，不要求两次自然语言说明逐字相同。方法收益由重算范围、处理时间和资源开销共同体现。

\section{构建上下文驱动的跨语言依赖恢复}
\label{sec:cross-language}

本项技术融合语法结构、编译语义和注册绑定规则，将不同语言中的声明、调用点和桥接入口连接为可追溯的工程关系。设计重点是明确关系的适用条件和推导依据，使检索能够区分实际确认的连接与尚待解析的候选。

\subsection{语法结构与编译语义融合}

Tree-sitter 提供多语言声明、引用和源码范围等结构事实，适合建立统一的基础表示。编译语义增强进一步利用类型、继承和符号表信息，确认在当前构建条件下可静态确定的引用目标。

构建上下文包含项目依赖、源集、编译选项和 SDK 等信息。例如，同一名称可能来自不同依赖版本，条件编译可能决定某个实现是否参与构建。语义结果必须关联对应上下文，才能在检索时解释其适用范围。

各语言分析前端将结果映射为统一的文件、符号、调用点和关系结构。支持程度按具体构建场景记录：具备语法事实时可提供结构查询，取得编译依据后再增强绑定。缺少 SDK 或无法静态确定的动态调用保留未解析状态，供后续工具或用户核查。

\subsection{跨语言注册绑定关系恢复}

JNI、N-API 等桥接机制通常通过注册名称、方法签名或函数引用关联上层接口和底层实现。关系恢复首先识别桥接入口，再解析注册表和绑定表达式，在对应模块与构建范围内匹配调用端和实现端。

以 Native 注册为例，上层调用使用的导出名称可在桥接声明中找到对应项，注册项进一步引用 C/C++ 实现。系统将两段已有依据的连接组合成跨语言路径，并保存调用位置、注册语句和实现声明作为证据。注册名相同但模块或签名不同的对象分别处理。

不同机制采用独立规则适配。KMP 的平台实现关联需要结合源集与声明约束，配置注入和事件关系需要结合注册键及使用位置。规则的输入、作用域和输出关系类型保持明确，便于使用可构建样例分别验证。

\subsection{统一依赖图与证据维护}

依赖图表示为 $G=(V,E)$。节点集合 $V$ 包含文件、符号及必要的绑定实体；边集合 $E$ 表达调用、继承、实现、注册绑定等关系。每条边同时关联来源版本、适用构建条件和证据集合，使图查询能够返回关系与其依据。

关系状态分为已解析、歧义和未解析。已解析表示在当前分析条件下找到有效目标；歧义保留多个候选；未解析记录缺少的端点或条件。证据来源与解析状态共同决定关系如何用于下游判断，不直接将名称相似度转换为确定调用。

更新时，系统从变化的源码或配置反查依赖这些证据的派生关系，撤回旧边并局部重算。新关系与相应源码版本发布后进入查询。影响分析通过有限邻域和路径查询展开，跨语言链路的质量按单边准确性与关键路径覆盖分别评价。

\section{层次路由与多视图融合的级联检索}
\label{sec:retrieval}

本项技术将职责层次、多粒度代码表示和混合召回结合。检索先用成本较低的索引形成候选，再对有限结果执行工程约束检查与精细比较，使复杂判断集中在更有价值的实现范围。

\subsection{显式粒度选择与自动路由}

系统提供自动、函数或方法、类或接口、模块、子系统五种检索粒度。用户指定粒度时，主要结果按该层次组织；已知代码位置作为定位锚点，不覆盖用户选择。必要接口、配置和关联实现仍可跨粒度补充。

自动模式依据任务、锚点和已发布索引组织候选，在判断不确定时保留多种粒度。粒度选择影响结果组织，不改变离线解析深度，也不要求重新建库。所选范围尚无对应索引时返回能力缺口，已经就绪但没有命中时返回空结果，两种情况分别处理。

任务描述还可包含目标语言、接口条件和平台约束。设计上将这些条件用于查询范围和候选筛选，避免将明确不适用的实现仅凭语义相关性提升到结果前列。

\subsection{层次检索与全局召回协同}

层次检索利用上层模块职责定位相关分支，再向下检索子模块、类和函数。上层摘要压缩了范围判断所需的信息，有助于减少进入细粒度比较的候选数量。

与此同时，全局通道直接查询函数与实现代码块，补充其他模块中的相关代码。两条路径共同形成候选集合：模块路径提供职责背景，全局路径降低模块边界或路由错误造成的漏检。跨模块任务可以同时命中不同分支的实现。

候选关联使用工程、版本和代码身份，避免同名对象被错误合并。函数命中可上溯所属模块取得职责信息，模块命中可继续读取局部实现。对比层次检索、全局检索及两者协同的结果，可以衡量范围收缩与实现召回之间的实际收益。

\subsection{多视图混合召回与排名融合}

代码表示按检索对象组织。函数与类保留声明、签名和实现片段；模块分别描述职责、公开接口和依赖。不同视图对应不同任务表述，使用户即使不知道具体标识，也能通过功能或接口条件进入候选范围。

全文检索匹配名称、签名和关键词，向量检索利用语义表示匹配自然语言与代码描述。HNSW 作为近似向量索引支撑候选召回，语义模型负责表示质量，两者分别影响检索效率和语义匹配能力。

多通道排名采用倒数排名融合，即 Reciprocal Rank Fusion（RRF）。设候选 $d$ 出现在若干检索通道中，$r_j(d)$ 表示其在通道 $j$ 中从 1 开始的名次，则融合分数为：
\begin{equation}
  S_{\mathrm{RRF}}(d)
  =\sum_{j:\,d\in L_j}\frac{1}{k+r_j(d)},\qquad k>0,
  \label{eq:rrf}
\end{equation}
其中 $L_j$ 为通道候选列表，$k$ 为平滑参数。未出现在某通道的候选不获得该通道贡献。该方法依据排名融合，减少全文与向量原始分数尺度不同带来的影响。

实际排序还需区分代码身份去重与视图信息合并。同一实现通过多个视图命中时，应合并为一个候选并保留其匹配依据；不同版本或不同工程的同名实现则保持独立。

\subsection{工程约束下的候选重排}

融合之后，系统对有限候选检查项目范围、版本、接口及平台条件，再进行更细的任务相关性比较。精排可以使用现有规则与可选模型，重点比较候选是否包含目标行为、接口是否适配以及依赖条件是否满足。

级联设计将大规模初筛交给索引，把高成本判断限制在较少候选。每阶段设置候选上限和时间预算，并将模型编码、重排和后续源码读取计入整体响应时间。超过预算时保留已得到的有效结果及缺口，不将未执行的阶段描述为已完成判断。

检索方法通过逐项对照评价：全文与符号基线、加入真实向量、加入模块表示、加入重排，分别观察定位质量、召回和尾延迟。跨语言代码相似性与具体目标环境的兼容性也分别判断，避免将功能描述相近直接解释为可迁移。

\section{任务覆盖与预算约束的上下文选择}
\label{sec:context-selection}

上下文构建以任务需要的证据为中心。输入是检索候选与依赖关系，输出是可以直接读取和核查的源码、接口及必要说明。核心设计包含覆盖建模、任务子图提取、预算化选择和缺口补查，形成从相关性到任务可用性的转换。

\subsection{任务需求与证据覆盖建模}

系统将任务分解为需要代码依据的要求，例如入口实现、接口约束、关键异常分支和配置来源。每份候选证据关联它支持的要求，多个片段可以共同支撑一项要求，同一片段也可能支撑多项要求。

为描述选择目标，设任务 $q$ 的必要要求集合为 $R(q)$，候选证据集合为 $\mathcal{E}_q$。令 $a_j(C)$ 表示证据子集 $C$ 是否满足要求 $j$ 的覆盖判定，$w_j$ 表示该要求的重要性，则设计目标可写为：
\begin{equation}
  C^{\star}=\arg\max_{C\subseteq\mathcal{E}_q}
  \sum_{j\in R(q)}w_j a_j(C),
  \label{eq:context-objective}
\end{equation}
满足
\begin{equation}
  \operatorname{Tokens}(\operatorname{Render}(C,q))\leq B,
  \quad |\operatorname{Files}(C)|\leq F,
  \quad \operatorname{Lines}(C)\leq L.
  \label{eq:context-budget}
\end{equation}

其中 $B$、$F$、$L$ 分别为 token、文件数和源码行数预算；$\operatorname{Render}$ 包含源码、来源、关系和缺口等最终交付文本。$a_j(C)$ 依据该要求的具体证据条件判定，例如实现正文与关键异常分支是否同时存在。覆盖权重和判定方式需要随任务集明确，不能由未经校准的模型置信度直接替代。

式~\eqref{eq:context-objective} 是覆盖驱动选择的目标描述，不代表当前已具备全局最优求解能力。在相近覆盖下减少 token 是进一步优化方向；实际效果仍需检查 Agent 是否能够完成任务。

\subsection{任务相关依赖子图提取}

依赖扩展从一个或多个候选实现出发，按任务确定方向与关系类型。接口变更优先关注调用方和实现类，迁移任务优先关注被依赖类型与配置，错误定位则围绕失败位置关联实现和异常路径。

扩展过程限制深度、节点数、高扇出邻居数量和时间，保留已访问集合以处理循环。对未解析关系记录缺口，对图上可达但与任务关联较弱的对象停止继续扩展。最终得到的任务子图为上下文选择提供候选，而不是将全部可达节点直接交付。

路径中的关系进一步映射到源码证据。调用记录用于解释连接，核心实现需要读取实际正文，接口约束可通过声明与签名表达，配置和测试按任务补入相应文本。多处证据共同解释一条跨语言关系时，应保留必要注册位置，避免只返回一条缺少依据的路径描述。

\subsection{预算约束的证据选择与源码去重}

选择过程优先保留主要实现及已确认的关键约束，再比较支持性证据对任务覆盖的新增贡献和文本成本。候选顺序不只取决于检索分数，还需要考虑证据角色：相似度较低的配置定义可能是解释行为所必需的信息。

源码去重首先依据版本、文件和范围识别相同或包含片段，再处理可合并的交叠范围。同一依赖被多个入口引用时只交付一次内容，保留必要关联说明。语义相似但属于不同实现或分支的代码不直接删除，避免将名称接近误判为内容等价。

最终预算以实际导出文本计量，计入标题、来源和缺口说明。单段证据超过剩余预算时，可以选择有明确边界的实现片段或保留后续读取线索；不能通过裁掉关键条件却仍宣称行为完整来满足预算。完整函数、局部片段和接口摘要需要在结果中有所区分。

当前上下文构建已采用核心实现优先、范围包含去重和精确 token 计量。覆盖贡献驱动选择、一般交叠范围合并及更完整的测试证据组织，按照上述目标继续增强。

\subsection{缺口识别与增量补取}

缺口主要来自三类情况：底层分析缺少必要关系，检索未找到足够证据，或预算不足导致部分内容未交付。系统保留缺口类型、相关代码位置和可继续查询的线索，支持 Agent 选择下一步读取对象。

连续查询记录 Agent 已持有证据的身份与内容哈希。相同内容不重复发送，变化内容重新提供。补取必须沿用可解释的来源版本；增强方案通过查询生命周期管理保护仍在使用的快照和结果。

当前 Agent 可利用基础符号、依赖和源码工具手动补查。进一步的自动补取将把未覆盖要求转换为局部查询，并使用剩余轮次、时间和 token 预算控制停止。整轮成本包含首包与后续补查，防止将上下文缩短的代价转移为更多工具调用。

\section{核心创新及方法协同}

项目的技术贡献集中在四项方法的协同。自适应分层建模将结构事实组织为职责层次，跨语言关系恢复补充功能实现的连接，多粒度检索在不同表述和粒度下定位候选，上下文选择则根据任务要求与预算确定最终交付内容。

\begin{table}[htbp]
  \centering
  \caption{核心方法与改进目标}
  \label{tab:technical-contributions}
  \small
  \begin{tabularx}{\textwidth}{P{0.20\textwidth}YY}
    \toprule
    核心方法 & 关键机制 & 改进目标 \\
    \midrule
    自适应分层并行建模 & 依赖约束划分、自适应终止、分支调度与局部增量 & 控制分析成本并形成有效职责结构 \\
    跨语言依赖恢复 & 构建语义融合、注册绑定规则、证据失效传播 & 提高依赖和关键调用路径覆盖 \\
    多粒度级联检索 & 层次与全局召回协同、多视图融合、有限候选重排 & 协调定位质量与在线时延 \\
    预算化上下文选择 & 任务子图、覆盖导向选择、范围去重与缺口补取 & 减少重复输入并保留必要信息 \\
    \bottomrule
  \end{tabularx}
\end{table}

模块树的质量影响范围路由，依赖图的质量影响证据覆盖，检索的候选质量影响上下文选择。因此，各层分别验证后，还需要进行组合对照，观察最终 Agent 的任务效果和总执行成本。基础解析器、索引算法与通信协议提供实现支撑，项目方法的价值通过上述组合机制及其测量结果体现。

% 第七章：当前产品实现与具体使用方式。
\chapter{产品实现与功能说明}
\label{chap:implementation}

\section{系统实现概况}

当前系统采用 TypeScript 单仓多包工程组织前后端及共享契约。代码索引服务负责源码扫描和结构解析，代码理解与检索服务负责模块产物、查询和上下文构建，适配服务负责分析与翻译。浏览器工作台和 VS Code 扩展复用主要交互组件，通过宿主与服务访问工程数据。

项目代码中仍保留 ForeXplore 的包名和命令前缀，本文统一以 RECAST 指称产品。SeekDB 保存版本化结构事实、模块产物和全文、向量检索文档，生成模型和向量模型分别通过服务配置接入。

表~\ref{tab:implementation-scope} 汇总当前产品能力。表中的“已接通”说明存在实际产品或服务通路，方法质量与业务收益需结合对应测试评价。

\begin{table}[htbp]
  \centering
  \caption{当前产品实现范围}
  \label{tab:implementation-scope}
  \small
  \begin{tabularx}{\textwidth}{P{0.19\textwidth}YY}
    \toprule
    功能 & 当前实现 & 继续增强方向 \\
    \midrule
    工程与结构索引 & 已接通快照、分批解析、持久化及局部查询 & 全流程资源约束与持续增量维护 \\
    模块层次 & 已接通父子模块、自适应状态及分页展示 & 真实模型职责判断与持久分支并行 \\
    任务检索 & 已接通全文、向量和多类文档候选查询 & 层次路由、工程约束与任务精排 \\
    依赖证据 & 已返回实际解析关系并有限补充源码 & 编译语义与完整跨语言绑定 \\
    上下文输出 & 已接通来源、范围去重、预算计量和导出 & 覆盖驱动选择与自动补取 \\
    Agent 接入 & 已接通 HTTP/MCP 任务与基础查询 & 端到端任务效果与连续反馈 \\
    复用迁移 & 已有多文件翻译及编译反馈后端首版 & 统一前端连接、审阅回填与行为验证 \\
    \bottomrule
  \end{tabularx}
\end{table}

\section{工程接入与代码建库}

\subsection{工程登记与项目选择}

用户通过工作台或扩展选择目标工程，按需登记参考工程。宿主为工程建立身份并识别其项目范围，索引服务记录当前源码版本。一个仓库内存在多个构建项目时，用户可选择本次关注的项目，浏览对应模块与源码。

工程身份、源码版本和本次任务角色分别管理。参考范围可以包含多个已登记工程，查询仅访问当前宿主可见的范围。用户输入路径用于工程登记，后续任务查询使用已发布的工程与项目身份。

\subsection{源码扫描与结构解析}

扫描阶段采集源码和相关工程文件，保存相对路径、文件角色、内容哈希及解析状态。构建元数据和测试文件分别标记，未支持或部分解析的文件保留在覆盖与诊断信息中。

结构解析采用可回收子进程处理批次，提取声明、源码范围和语法依赖。当前大仓解析批次最多包含 64 个文件或 8 MiB 输入。结构事实绑定后分批写入数据库，避免在写入阶段逐条处理大量记录。

解析结果用于支持后续模块建模和源码查询。界面分别展示文件归属、解析情况和依赖缺口，使用户能够了解结果依据；这些维度不会合并为一个“理解完成率”。

\subsection{索引发布与模块重建}

基础索引在构建版本完成后发布，查询读取已就绪的源码和结构。模块理解结果及其检索投影在此基础上更新，基础源码查询可以先于完整模块说明使用。

当前支持复用已有源码索引重建模块树与对应投影，无需再次执行基础解析。旧模块产物保持兼容，树和投影状态按当前范围显示。历史源码快照可以用于只读查询，但历史模块产物是否可检索还取决于相应投影是否保留且可用。

\section{模块建模与层次展示}

\subsection{自适应模块产物}

当前模块产物通过父子关系表达层次，节点记录模块或子系统类型、范围规模、细化状态及判断原因。结构判断依据目录、接口和依赖形成划分，可选模型规划器在有限证据上提出拆分或停止决定。

系统区分已拆分、有效叶节点和暂未细化的范围。职责足够单一的小工程可以直接形成叶模块；多职责范围继续细分；证据或预算不足的范围保留原因。候选组数量上限用于控制单次判断输入，与解析批次和模块文件规模是不同参数。

目前真实工作台的层次验证使用本机结构判断。模型决策接口与校验逻辑已接入，实际采用模型还需配置可用规划器；真实业务命名和边界质量应在模型运行之后独立评价。

\subsection{父子范围与模块导航}

只有终端模块直接保存文件归属，父节点通过子树汇总范围。宿主检查树结构无环、父节点存在、子组互斥和文件覆盖，模块统计按聚合范围计算，避免父子重复计数。

前端展示根模块、子模块及其目录、文件和符号。节点按页读取，首屏只发送有限层次与全量统计，用户展开后继续取得子节点。搜索覆盖宿主已索引的完整节点范围，未展开符号也可被定位。

选择模块后，详情区显示职责、接口、依赖、细化原因和判断来源。模块层次表示职责范围，目录层次表示原始源码位置，两者在导航中分别保留。

\section{自然语言与多粒度任务检索}

任务检索页提供需求输入、工程范围、粒度选择和内容预算。用户可以直接描述任务，也可以先在模块树中限定范围。主要粒度包括自动、函数或方法、类或接口、模块及子系统。

当前检索查询声明、实现代码块和模块摘要，结合全文与向量召回组织候选。自动模式根据实际可用候选形成多粒度结果；手动模式按选择组织主要实现，支持性依赖可跨粒度提供。模型意图规划和更完整的层次路由作为增强路径继续接入。

子系统选项按产物和投影能力启用。当前范围只有普通模块，或子系统检索投影尚未就绪时，界面显示不可用，接口返回对应缺口。函数和类粒度同样依据范围内的索引能力判断，不通过静默替换改变用户选择。

语义检索需要配置与索引向量一致的 embedding 模型。已有中文查询使用 multilingual E5 进行验证；未配置真实模型时的确定性 hash 向量用于容量和接口测试，不作为自然语言语义质量依据。更换模型或维度需要同步重建相应向量空间。

\section{依赖关系与源码证据}

当前任务结果返回实际解析的依赖记录，包含来源符号或文件、目标引用、关系类型和解析状态。服务在有限范围内补充相关声明和源码，使用户可以从候选入口继续理解其关联实现。

源码预览读取固定快照，并携带行列范围与内容哈希。用户能够核对命中的实现正文，以及引用关系是否有相应代码依据。未解析端点保留状态，跨语言桥接能力随对应语义与规则分析逐步增强。

产品当前提供的有限依赖证据与第~\ref{sec:cross-language} 节的完整图谱方案具有不同范围。完整多跳路径、跨语言调用与影响分析需具备相应关系后使用，现有结构查询作为基础入口保留。

\section{任务上下文生成与导出}

\subsection{输出内容与预算计量}

上下文构建器从候选中优先读取主要实现，再补入有限支持性证据。返回包包含任务、工程快照、粒度结果、相关实现、源码证据、关系、缺口和用量，文本输出由同一份内容生成。

当前实现对最终完整 Markdown 使用 \texttt{cl100k\_base} 计量 token，同时执行文件数和源码行数约束。计量范围包括源码及来源、关系、缺口等文本，避免仅统计源码正文而低估实际交付量。

已知证据通过身份和内容哈希去重，范围包含时避免重复发送相同源码。预算不足或分析存在缺口时返回部分结果及原因。这里的完成状态描述本次查询流程，不代表已证明覆盖全部业务行为。

\subsection{结果预览与内容选择}

界面提供证据列表、源码预览、选择、复制和下载。默认导出使用服务端已计量的文本，显示对应 token 用量。用户手动筛选证据后，界面区分原包用量与新内容的待计量状态，防止沿用已经不对应当前文本的数值。

需求、范围、粒度或预算变化会撤销旧请求，结果按请求身份更新。取消操作向服务传播，已失效请求的迟到结果不会覆盖新任务。

\section{Coding Agent 接入}

\subsection{HTTP 接口}

任务 HTTP 接口接收需求、固定工程范围、检索粒度和预算，返回完整结构化上下文包。主要字段按表~\ref{tab:task-contract} 组织，客户端无需了解数据库表或内部模块实现。

\begin{table}[htbp]
  \centering
  \caption{任务检索的主要输入与输出}
  \label{tab:task-contract}
  \small
  \begin{tabularx}{\textwidth}{P{0.16\textwidth}YY}
    \toprule
    方向 & 内容 & 用途 \\
    \midrule
    请求 & 任务描述、工程与版本、粒度、预算 & 确定查询范围和结果组织 \\
    请求 & 已持有证据及内容哈希 & 避免重复源码输入 \\
    响应 & 实现候选、实际粒度、来源快照 & 明确定位结果和适用范围 \\
    响应 & 源码、关系、范围与哈希 & 提供可核验的开发依据 \\
    响应 & 内容用量、部分结果与缺口 & 决定是否继续补查 \\
    \bottomrule
  \end{tabularx}
\end{table}

\subsection{MCP 工具}

MCP 通过 \texttt{search\_task\_context} 提供任务上下文查询，并复用工程列表、项目列表、符号查询、依赖查询和源码读取工具。Agent 先取得允许访问的工程及版本，再发起任务查询；需要更多信息时使用基础工具局部补充。

任务工具输出一份已计量的 Markdown，避免同时发送包含相同源码的多份表示。当前适配器连接本机查询服务，工程登记、文件访问和范围检查由宿主完成。将接口提供给外部 Agent 与验证其完整任务效果分别开展。

\subsection{IDE 宿主与浏览器工作台}

VS Code 扩展负责当前工程选择、模块导航、任务输入和结果展示，支持从源码位置进入后续开发。浏览器工作台复用主要界面与服务通路，用于独立浏览、查询和演示。

界面顶部保留任务检索与复用迁移两个入口，选择目标、输入需求和选择方案占据主要操作区；模块说明、覆盖统计和解析诊断放在下方。信息层次围绕常用操作组织，使开发者能够从候选回到源码，再返回当前任务继续筛选。

\section{复用迁移与翻译后端}

当前翻译后端接收上下文与任务规格，先由分析阶段确定修改计划，再执行多文件生成及编译反馈修复。运行记录保存计划、变更和编译结果，并提供任务查询、取消、恢复与回退能力。

任务检索界面可以携带需求进入复用迁移流程，但选定证据直接接入完整多文件翻译、审阅与回填仍需统一。已有独立迁移宿主提供补丁验证和审阅机制，后续需要将生成阶段的产物与其接通，而不是仅在界面上增加流程步骤。

编译通过说明目标语言和工程检查得到满足。迁移是否保留异常处理、资源释放及其他必须行为，需要依据目标任务补充行为验证。翻译输入、修改记录和验证结果应对应同一任务与目标基线。

\section{运行方式与典型使用流程}

\subsection{运行环境}

浏览器工作台需要 Node.js 24、已安装的工程依赖、SeekDB，以及与数据库向量维度匹配的 embedding 服务。配置生成模型后可启用相应模块理解和翻译能力。目标语言编译器按实际迁移任务准备。

本地工作台从项目根目录启动，指定目标工程及可选参考工程：
\begin{verbatim}
npm run dev:code-workbench -- --target /path/to/project
\end{verbatim}
服务输出实际页面和查询地址。已有源码索引时可选择重新构建模块产物，模型和数据库参数按项目运行配置设置。凭据保留在宿主或服务环境中。

\subsection{典型任务操作}

以文件上传限制修改为例，用户先选择已登记的 FileUpload 工程，在任务检索页输入需求，指定函数粒度及内容预算。系统返回相关实现和源码，用户核对总大小、单文件限制与异常处理分支，再查看依赖和缺口。

确认结果后，用户可以复制或下载上下文，也可以由 Agent 通过 MCP 取得同一任务结果。若仍需了解某个类型或配置，Agent 按返回位置继续查询。选择复用迁移时，用户进一步指定参考实现与目标约束，进入分析和生成流程。

该使用流程连接了“需求描述、实现定位、证据核查、上下文交付”四个实际环节。持续开发中的修改验证、索引更新和定向补查在相同工程身份与任务约束下继续衔接。

% 第八章按要求保留结构，正文留空。
\chapter{项目进展与阶段性验证}
\label{chap:validation}
\section{项目研发起点与开发时间线}
\section{版本迭代与当前完成情况}
\subsection{已完成功能}
\subsection{正在开发功能}
\subsection{后续规划功能}
\section{测试数据、环境与评价方法}
\section{阶段性验证结果}
\subsection{建库规模、资源开销与增量更新效率}
\subsection{模块划分质量与自适应深度}
\subsection{跨语言关系准确率与路径覆盖}
\subsection{检索质量与查询时延}
\subsection{上下文证据覆盖与 token 开销}
\subsection{Coding Agent 任务效果对照}
\section{关键方法消融实验}
\section{命题指标对应关系}
\section{下一阶段开发与验证计划}
\subsection{千万行级工程与企业场景验证}
\subsection{任务集扩展与持续回归}

% 第九章按要求保留结构，正文留空。
\chapter{项目团队与组织分工}
\section{团队总体情况与组织结构}
\section{成员专业背景与能力构成}
\section{团队成员任务分工}
\subsection{超大工程索引与模块建模}
\subsection{跨语言依赖分析}
\subsection{智能检索与上下文构建}
\subsection{系统后端与接口开发}
\subsection{产品前端与交互设计}
\subsection{测试与实验评估}
\section{成员实质性贡献}
\section{指导教师分工}
\section{师生协同机制}
\section{企业、教师与学生协同机制}
\section{项目组织与研发管理}

% 第十章按要求保留结构，正文留空。
\chapter{项目总结与后续规划}
\section{项目阶段性总结}
\section{当前成果与主要不足}
\section{核心技术迭代规划}
\section{产品落地与集成规划}
\section{规模化验证规划}
\section{企业应用与推广规划}

\end{document}
